%%%%%%%%%%%%%%%%%%%%%%%%%%%%%%%%%%%%%%%%%%%%%%%%%%%
%
%  CHAPTER I -- INTRODUCTION AND BACKGROUND
%
%%%%%%%%%%%%%%%%%%%%%%%%%%%%%%%%%%%%%%%%%%%%%%%%%%%

\pagestyle{plain} % No headers, just page numbers
\pagenumbering{arabic} % Arabic numerals
\setcounter{page}{1}

\chapter{\uppercase{Introduction and Background}}
\label{cha:introduction}

\section{Fibrillar Structures: The Whole Is Greater Than the Sum of Its Parts}
\label{sec:motivation}

A woven fabric is made of threads, but none of its properties belong to a single thread. Whether the fabric hangs together, how it drapes, and how stiff it is in each direction all follow from how the threads are interlaced. The same holds for knitted garments, where one continuous yarn becomes a stretchable surface, and for chainmail, where rigid rings become a flexible sheet. In each case the whole is greater than the sum of its parts.

We use the term \emph{fibrillar structures} for assemblies of this kind: collections of one-dimensional elements (threads, strands, rings, or spatial graphs) whose global behavior comes from their arrangement. They appear at every scale. Woven and three-dimensional (3D) woven composites carry loads in aerospace and automotive parts~\cite{mouritz1999review,STIG20121495}. Chainmail-inspired and polycatenated architected materials switch between fluid-like and solid-like responses~\cite{tang2023chainmail,zhou20253d}. Knotted and intertwined microarchitectures improve toughness~\cite{moestopo2023knots}. In chemistry and biology, catenated molecules and fibrillar tissues show the same principle at the molecular scale~\cite{carlucci2003polycatenation} (Figure~\ref{fig:motivation-montage}).

\begin{figure}[htbp]
\centering
\includegraphics[width=5.5in]{ch01_fig01_motivation_montage.png}
\caption[Fibrillar structures across scales.]{Fibrillar structures across scales: (a) a woven textile, (b) a woven composite, (c) a knitted structure, and (d) chainmail. In each case the assembly has properties that none of its threads or rings has on its own.
\figtodo{Compose four panels. (a), (b) Prelim\_v2 slide~5 (``Motivation'': woven carpet, woven composites; media \fs{image7.png}, \fs{image8.jpeg}--\fs{image11.png}). (c) a knit: \cite{yildiz2026hypergraph} Fig.~\fs{fig:bravais2D_repeated_nonuniform} panel \fs{images/ws_2d/2_2_1.jpg}, or a photo. (d) chainmail: \cite{yildiz2026hypergraph} Fig.~\fs{fig:tetra_link} (\fs{images/tetchain/0.png}). Slide-5 photos are third-party: add credits to the caption.}}
\label{fig:motivation-montage}
\end{figure}

\section{The Design Problem}
\label{sec:challenge}

If the behavior of a fibrillar structure lives in its arrangement, then designing one means designing an arrangement. This requires a representation that separates local choices from global structure, so that a designer can make a small decision at one place and know what it does to the whole.

For planar weaves such a representation has existed for decades~\cite{Grunbaum80}. For volumetric weaves, links, knits, and catenanes it largely does not. Volumetric weaves are usually designed by stacking planar weaves and adding straight through-thickness threads~\cite{STIG20121495}. Linked and polycatenated materials are typically found by inspection or taken from isolated chemical examples~\cite{zhou20253d}. Knitted and entangled materials are often generated by simulating the physical process that produces them~\cite{moestopo2023knots}. Each approach works within one family. None of them describes a weave, a knit, and a chainmail with the same language, and none gives a design space that can be counted, enumerated, or explored by rule.

This dissertation addresses that gap. Its premise is that weaves, knits, and chainmail are not separate design problems but instances of one.

\section{Background}
\label{sec:related-work}

We review four lines of work that the remainder of this proposal builds on.

\subsection{Planar Weaves and Topological Labeling}
Gr\"unbaum and Shephard classified planar isonemal fabrics by encoding each crossing with a binary over/under label and studying the symmetry group of the resulting pattern~\cite{Grunbaum80,Grunbaum85,Grunbaum86,grunbaum1988isonemal}. A weave becomes a labeling of a grid, and questions about the weave become questions about the labeling. The same idea of a discrete local label appears in later work on stitch meshes for knitwear~\cite{yuksel2012stitch}, interactive 3D weave design~\cite{wu2020weavecraft}, and interlocking tiles that realize bi-axial weaves~\cite{krishnamurthy2020bi,krishnamurthy2021geometrically}. Knitted fabrics have also been described through their geometry and through bicontinuous surfaces~\cite{knittel2020modelling,niu2025geometric}.

\subsection{Rotation Systems and Cyclic Weaving}
A graph alone does not determine a set of cycles~\cite{tait1877knots}. A rotation system, the cyclic order of edges around each vertex, does~\cite{edmonds1960combinatorial,gross2001topological}. Rotation systems are implicit in winged-edge and half-edge mesh structures, and Akleman and colleagues used them to weave arbitrary 2-manifold meshes: assigning a twist to each edge turns face boundaries into interlaced strands, giving plain, twill, and single-cycle weaves on surfaces~\cite{akleman2009cyclic,xing2010single,akleman2011cyclic,akleman2015extended,akleman2020topologically}. These methods assume a 2-manifold surface and binary twists. Block meshes extend the setting to graphs embedded in 3-manifolds~\cite{akleman2015blockmeshes}, but knots and links appear there only implicitly.

\subsection{Volumetric Weaves and Periodic Entanglements}
Volumetric weaves are related to rod packings, whose symmetries are well catalogued~\cite{okeeffe2001invariant_cubic_rod_packing,rosi2005rodpacking}, and to periodic entanglements obtained by mapping hyperbolic line patterns onto triply periodic surfaces~\cite{evans2013periodiceentanglement1,evans2013periodiceentanglement2}. Both describe a structure through one global embedding. Neither assigns independent local configurations to the cells of a lattice.

\subsection{Interlocked and Architected Materials}
Mechanically interlocked architectures include links, catenanes, and chainmail, where components are inseparable without being fused~\cite{carlucci2003polycatenation,tang2023chainmail,zhou20253d}. On the materials side, triply periodic minimal surfaces (TPMS) and plate lattices are the standard families of architected cellular solids~\cite{schoen1970infinite,alketan2019tpms,berger2017plate}, and their mechanical response is usually interpreted through the scaling laws of cellular solids~\cite{gibson1997cellular}. These serve as the baselines for the mechanical study proposed in Chapter~\ref{cha:mechanical}.

\section{Thesis Statement}
\label{sec:goal}

\begin{quote}
\emph{Woven, knitted, linked, and catenated structures can be designed within one framework, in which a scaffold carries a small set of local labels and a replacement rule turns each labeled element into fiber. Hypergraph rotation systems realize this framework for all cycle-based fibrillar structures with a single local operator, and structures designed in it can be manufactured and measured to show how arrangement alone determines physical behavior.}
\end{quote}

Table~\ref{tab:scheme} lists the constructions developed in this dissertation in terms of these three ingredients. Each row was a separate paper with its own construction. Read together, they show the same pattern: the scaffold grows more general, the labels stay small, and the replacement rule converges to the single twist operator of a hypergraph rotation system.

\begin{table}[htbp]
\centering
\begin{singlespace}
\footnotesize
\begin{tabular}{@{}>{\raggedright\arraybackslash}p{0.17\textwidth}>{\raggedright\arraybackslash}p{0.18\textwidth}>{\raggedright\arraybackslash}p{0.16\textwidth}>{\raggedright\arraybackslash}p{0.22\textwidth}>{\raggedright\arraybackslash}p{0.17\textwidth}@{}}
\toprule
\textbf{Construction} & \textbf{Scaffold} & \textbf{Local label} & \textbf{Replacement rule} & \textbf{Produces} \\
\midrule
Woven tiles (Ch.~\ref{cha:modular-planar}) & Square grid & Over/under bit & Cell $\to$ congruent tile & 2-way 2-fold weaves \\[3pt]
Volumetric Bravais weaves (Ch.~\ref{cha:volumetric-bravais}) & Bravais lattice, Wigner-Seitz cells & Skew-line configuration & Cell $\to$ thread segments & Volumetric weaves, rod packings \\[3pt]
Volumetric nonwovens (Ch.~\ref{cha:nonwoven}) & Cubic grid & Octal connection digit & Cell $\to$ attachments & Frames, infinite polyhedra \\[3pt]
Linked structures (Ch.~\ref{cha:linked}) & Mesh or Wigner-Seitz honeycomb & Edge bundled or not & Face $\to$ face-cycle, edge $\to$ edge-cycle & Links, chainmail \\[3pt]
Primal-dual polycatenanes (Ch.~\ref{cha:linked}) & Labeled 3-honeycomb and dual & Cell selected or not & Cell $\to$ spatial graph & Polycatenanes \\[3pt]
\textbf{Hypergraph rotation systems (Ch.~\ref{cha:hgrs})} & \textbf{Graph with faces, finite or periodic} & \textbf{Integer twist per edge} & \textbf{Face-sides reconnected by a cyclic shift} & \textbf{Knots, links, weaves, knits, chainmail} \\
\bottomrule
\end{tabular}
\end{singlespace}
\caption[Constructions of this dissertation as scaffold, label, and rule.]{The constructions of this dissertation described by the same three ingredients. The last row unifies the cycle-based cases above it (Section~\ref{sec:hgrs-synthesis}).}
\label{tab:scheme}
\end{table}

\section{Contributions}
\label{sec:contributions}

The completed and proposed contributions are:

\begin{enumerate}
\item[\textbf{C1}] \textbf{Modular planar weaves.} Two congruent, foldable tiles that realize every 2-way 2-fold weave by translation and rotation~\cite{yildiz2023modular}.
\item[\textbf{C2}] \textbf{Volumetric Bravais weaves.} A definition of volumetric weaves through the isometries of Bravais lattices and Wigner-Seitz cells, with design strategies, lower bounds, and an enumeration of cubic primitive (cP) weaves~\cite{yildiz2025volumetric,akpanya2025enumeration}.
\item[\textbf{C3}] \textbf{An algebra for labeled lattices.} A numeral system of $P^K$-adic zeros whose operations generate periodic label arrays, introduced for volumetric nonwovens~\cite{yildiz2024nonwoven}.
\item[\textbf{C4}] \textbf{Linked and catenated structures.} Triply periodic linked structures from labeled meshes~\cite{yildiz2025linked} and primal-dual polycatenanes from labeled 3-honeycombs~\cite{yildiz2026polycatenane}.
\item[\textbf{C5}] \textbf{Hypergraph rotation systems.} A combinatorial representation, based on cyclic edge-face incidences and integer edge twists, that unifies knots, links, weaves, knits, and chainmail~\cite{yildiz2026hypergraph}. This is the unifying contribution of the dissertation.
\item[\textbf{C6}] \textbf{Manufacturing volumetric weaves} (proposed). A set of fabrication methods for volumetric weaves, building on computational fabrication of optical channels~\cite{deshpande2025optical}, compared on what part of the design space each can realize.
\item[\textbf{C7}] \textbf{Mechanical exploration of cP-weave lattices} (proposed). A compression study of cP-weave lattices against TPMS baselines at a fixed volume fraction, which isolates the effect of arrangement on mechanical response.
\end{enumerate}

\section{Organization}
\label{sec:organization}

Figure~\ref{fig:roadmap} shows how the chapters connect. Chapters~\ref{cha:modular-planar} to~\ref{cha:linked} develop the scheme of Table~\ref{tab:scheme} one construction at a time, and each ends with the lesson that motivated the next. Chapter~\ref{cha:hgrs} presents hypergraph rotation systems and shows how they unify the earlier constructions, and where they do not. Chapters~\ref{cha:manufacturing} and~\ref{cha:mechanical} present the proposed work on manufacturing and mechanical exploration. Chapter~\ref{cha:summary} summarizes the contributions and gives a timeline.

\begin{figure}[htbp]
\centering
\includegraphics[width=6in]{ch01_fig02_roadmap.png}
\caption[Roadmap of the dissertation.]{Roadmap of the dissertation. The design chapters move from a planar grid to lattices, an algebra on labels, labeled meshes and honeycombs, and finally hypergraph rotation systems. The proposed work connects these designs to fabrication and measurement.
\figtodo{New schematic. One thumbnail per chapter, left to right: woven tiles (\cite{yildiz2023modular} Fig.~\fs{fig:planarslabs}, \fs{2D/results/slab1.png}), cP-weave (\cite{yildiz2025volumetric} \fs{fig1-teaser.pdf}), nonwoven frame (\cite{yildiz2024nonwoven} \fs{images/flexible_wood/wood3.png}), linked structure (\cite{yildiz2025linked} \fs{images/ws/CP.jpg}), H-GRS (\cite{yildiz2026hypergraph} \fs{images/teaser/cube/2man_chainmail.jpg}); then a second row for fabrication (Prelim\_v2 slide~56) and mechanical testing (Prelim\_v2 slide~71). Arrows labeled with each chapter's ``lesson''.}}
\label{fig:roadmap}
\end{figure}
